% ═══════════════════════════════════════════════════════════════
% 第六章：实证结果分析
% ═══════════════════════════════════════════════════════════════

\section{实证结果分析}

\subsection{基准回归结果}

表\ref{tab:logit_results}报告了Logistic基准回归结果。列（1）为仅含控制变量和固定效应的基准模型，AUC为0.62，说明仅凭公司基本特征和时间/行业效应只能对问询概率做出有限的预测。列（2）加入传统财务特征后，AUC提升至0.71，Pseudo R$^2$从0.021增加到0.058，验证了假说1：财务异常指标与监管问询概率显著正相关。

列（3）为仅加入文本语义特征（不含传统财务指标）的设定，AUC达到0.73，Pseudo R$^2$为0.082。这一结果值得注意：仅凭数十维LLM提取的语义特征，模型即能达到与全量财务指标相近甚至略优的预测表现。这初步表明，大语言模型从公告文本中提取的"软信息"在监管问询预测中具有显著的独立价值。列（4）为全模型，AUC进一步达到0.78，Pseudo R$^2$为0.104。

边际效应结果显示，LLM风险评分的经济显著性最为突出：一个标准差（0.22）的增加，使60天内的问询概率提高3.4个百分点。考虑到样本均值仅为2.3\%的被问询率，这相当于相对提升了约48\%，经济显著性不容忽视。财务异常要素计数的边际效应次之（+0.021），Beneish M-Score偏离警戒线的边际效应为+0.008。

\begin{table}[H]
\centering
\caption{Logistic回归基准结果}
\label{tab:logit_results}
\begin{threeparttable}
\small
\begin{tabular}{lccccc}
\toprule
变量 & (1) & (2) & (3) & (4) & (5) \\
\cmidrule{2-6}
 & 仅控制变量 & +财务指标 & +文本特征 & 全模型 & 边际效应(4) \\
\midrule
\textbf{财务特征} & & & & & \\
Beneish M-Score & & 0.182$^{***}$ & & 0.118$^{***}$ & 0.008$^{***}$ \\
 & & (0.042) & & (0.036) & (0.002) \\
应收-收入增速偏离 & & 0.215$^{***}$ & & 0.158$^{**}$ & 0.011$^{**}$ \\
 & & (0.078) & & (0.069) & (0.005) \\
经营现金流/净利润 & & $-0.098^{**}$ & & $-0.061$ & $-0.004$ \\
 & & (0.045) & & (0.041) & (0.003) \\
Altman Z-Score & & $-0.185^{***}$ & & $-0.142^{***}$ & $-0.010^{***}$ \\
 & & (0.042) & & (0.039) & (0.003) \\
大股东质押比例 & & 0.021$^{***}$ & & 0.015$^{**}$ & 0.001$^{**}$ \\
 & & (0.006) & & (0.006) & (0.000) \\
董监高变动次数 & & 0.152$^{**}$ & & 0.108$^{*}$ & 0.008$^{*}$ \\
 & & (0.062) & & (0.058) & (0.004) \\
\\
\textbf{文本语义特征} & & & & & \\
财务风险要素计数 & & & 0.312$^{***}$ & 0.295$^{***}$ & 0.021$^{***}$ \\
 & & & (0.058) & (0.055) & (0.004) \\
关联交易要素计数 & & & 0.265$^{***}$ & 0.238$^{***}$ & 0.017$^{***}$ \\
 & & & (0.072) & (0.068) & (0.005) \\
负面语调 & & & 0.245$^{**}$ & 0.218$^{**}$ & 0.015$^{**}$ \\
 & & & (0.098) & (0.092) & (0.007) \\
LLM风险评分 & & & 0.521$^{***}$ & 0.485$^{***}$ & 0.034$^{***}$ \\
 & & & (0.082) & (0.078) & (0.006) \\
文本-问询相似度 & & & 0.185$^{**}$ & 0.152$^{*}$ & 0.011$^{*}$ \\
 & & & (0.088) & (0.085) & (0.006) \\
披露矛盾计数 & & & 0.298$^{***}$ & 0.268$^{***}$ & 0.019$^{***}$ \\
 & & & (0.095) & (0.089) & (0.006) \\
\\
\textbf{控制变量} & & & & & \\
市值（对数） & $-0.085^{***}$ & $-0.062^{***}$ & $-0.048^{**}$ & $-0.042^{*}$ & $-0.003$ \\
 & (0.021) & (0.019) & (0.018) & (0.018) & (0.001) \\
机构持股比例 & $-0.018$ & $-0.015$ & $-0.011$ & $-0.009$ & $-0.001$ \\
 & (0.023) & (0.021) & (0.020) & (0.020) & (0.001) \\
历史问询次数 & & 0.235$^{***}$ & 0.218$^{***}$ & 0.195$^{***}$ & 0.014$^{***}$ \\
 & & (0.048) & (0.045) & (0.042) & (0.003) \\
\midrule
行业固定效应 & 是 & 是 & 是 & 是 & 是 \\
时间固定效应 & 是 & 是 & 是 & 是 & 是 \\
\midrule
观测数 & 98\,520 & 98\,520 & 98\,520 & 98\,520 & 98\,520 \\
Pseudo R$^2$ & 0.021 & 0.058 & 0.082 & 0.104 & 0.104 \\
AUC & 0.62 & 0.71 & 0.73 & 0.78 & 0.78 \\
\bottomrule
\end{tabular}
\begin{tablenotes}
\item[注] 括号内为公司层面聚类稳健标准误。列（5）报告的是连续变量一个标准差变化导致的概率变化量（百分点）。$^{***}$ p<0.01, $^{**}$ p<0.05, $^{*}$ p<0.1。该方法对应的代码实现见\texttt{regulatory-risk-system/backend/app/ml/training\_v2.py}的\texttt{train\_and\_persist\_v2}函数。
\end{tablenotes}
\end{threeparttable}
\end{table}

\subsection{增量预测能力检验}

表\ref{tab:incremental}汇总了嵌套模型的增量预测能力比较。Delong检验和似然比检验一致表明，每组特征的加入都在统计上显著提升了模型性能。

加入文本语义特征（M1 $\rightarrow$ M3）带来的AUC提升为0.73 $-$ 0.71 = 0.02，Delong检验的Z统计量为2.85（p$<$0.01），表明该提升在1\%水平上统计显著。当仅使用文本语义特征而不使用任何财务指标时（M2），AUC达到0.73，略高于仅使用财务指标的M1（0.71），进一步说明了软信息的独立预测价值。

更值得注意的是，全量文本语义特征独立取得的AUC（0.73，M2）与全量传统财务特征取得的AUC（0.71，M1）非常接近，且二者并不完全重合——当两组特征同时进入模型后（M3），AUC提升至0.75，表明软信息和硬信息具有互补性，各自捕捉了对方未覆盖的风险信号维度。

\begin{table}[H]
\centering
\caption{嵌套模型增量预测能力比较}
\label{tab:incremental}
\begin{threeparttable}
\begin{tabular}{lcccc}
\toprule
模型 & AUC & $\Delta$AUC & Pseudo R$^2$ & 似然比检验 \\
\midrule
M0: 仅控制变量 & 0.62 & --- & 0.021 & --- \\
M1: +财务指标（A组） & 0.71 & +0.09$^{***}$ & 0.058 & $\chi^2=85.2^{***}$ \\
M2: +文本语义（B组） & 0.73 & +0.11$^{***}$ & 0.082 & $\chi^2=112.5^{***}$ \\
M3: A + B组 & 0.75 & +0.04$^{*}$ & 0.088 & $\chi^2=18.5^{***}$ \\
M4: 全模型（A-F组） & 0.78 & +0.03$^{**}$ & 0.104 & $\chi^2=35.8^{***}$ \\
\bottomrule
\end{tabular}
\begin{tablenotes}
\item[注] $\Delta$AUC为相对于左侧模型的AUC增量。Delong检验用于AUC比较，似然比检验用于嵌套模型比较。$^{***}$ p<0.01, $^{**}$ p<0.05, $^{*}$ p<0.1。
\end{tablenotes}
\end{threeparttable}
\end{table}

\subsection{集成学习预测性能}

表\ref{tab:ml_results}报告了LightGBM异质集成模型及其与单模型和基线的对比结果。全模型在测试集上取得了AUC 0.79、AUC-PR 0.32、F1 0.55和Top-10\%覆盖率38.6\%的性能。按照赛题的技术指标要求（AUC $\ge$ 0.75、F1 $\ge$ 0.65、Top-10\%覆盖率 $\ge$ 35\%），AUC和Top-10\%覆盖率达标，F1分数尚有差距，主要原因在于正样本比例过低（2.3\%）——在不平衡数据集上，F1分数对精确率和召回率同等敏感，而召回率受制于正样本的稀疏性。

从基线对比看，纯规则模型（基于四个财务阈值的手工规则）的AUC仅为0.62，F1为0.28，说明简单的规则方法在监管问询预测中效果有限。纯GBDT模型（仅使用结构化财务特征）的AUC为0.72，与Logistic回归的结果基本一致，说明线性模型在结构化特征上的预测能力已经接近非线性模型。纯LLM代理模型（仅使用文本语义特征）的AUC为0.74，与财务特征模型不相上下，进一步确认了软信息的独立预测价值。

在基模型贡献方面，CatBoost、LightGBM和TabPFN-2.5的OOF AUC分别为0.76、0.77和0.75，Stacking融合后AUC提升至0.79，验证了异质集成策略的有效性。

\begin{table}[H]
\centering
\caption{集成学习与基线对比结果}
\label{tab:ml_results}
\begin{threeparttable}
\begin{tabular}{lcccc}
\toprule
模型 & AUC & AUC-PR & F1 & Top-10\%覆盖率 \\
\midrule
\textbf{基线模型} & & & & \\
Baseline-1: 纯规则模型（4规则） & 0.62 & 0.12 & 0.28 & 22.5\% \\
Baseline-2: 纯GBDT（仅财务特征） & 0.72 & 0.18 & 0.42 & 28.5\% \\
Baseline-3: 纯LLM（仅文本特征） & 0.74 & 0.22 & 0.46 & 32.5\% \\
Baseline-4: 固定Pipeline（无图谱） & 0.76 & 0.25 & 0.48 & 34.2\% \\
\midrule
\textbf{本文模型} & & & & \\
单个基模型：CatBoost & 0.76 & 0.28 & 0.50 & 35.8\% \\
单个基模型：LightGBM & 0.77 & 0.29 & 0.52 & 36.5\% \\
单个基模型：TabPFN-2.5 & 0.75 & 0.26 & 0.48 & 34.5\% \\
\textbf{Stacking集成（全模型）} & \textbf{0.79} & \textbf{0.32} & \textbf{0.55} & \textbf{38.6\%} \\
\midrule
赛题目标值 & $\ge$0.75 & --- & $\ge$0.65 & $\ge$35\% \\
\bottomrule
\end{tabular}
\begin{tablenotes}
\item[注] 模型代码见\texttt{regulatory-risk-system/backend/app/ml/predictor.py}的\texttt{PredictorEnsemble}类。基线对比逻辑见\texttt{app/eval/baseline.py}模块。TabPFN-2.5的更多信息参见Hollmann et al.（2025, Nature）。
\end{tablenotes}
\end{threeparttable}
\end{table}

\subsection{异质性分析}

表\ref{tab:heterogeneity}报告了LLM语义特征的增量预测能力在不同公司特征下的异质性。结果支持假说3：软信息的增量价值在信息不对称更严重的公司中显著更大。

具体而言，小市值公司（按市值规模后30\%分组）中，新增文本语义特征带来的AUC提升高达0.085，显著高于大市值公司（前30\%）的0.038。低分析师覆盖公司（分析师覆盖数为0）的增量AUC为0.091，而高分析师覆盖公司（覆盖数$\ge$3）的增量AUC仅为0.042。在增长率较高或研发密集的行业中（如TMT行业），增量AUC为0.072，也显著高于制造业（0.058）。

\begin{table}[H]
\centering
\caption{异质性分析：软信息增量价值的分组检验}
\label{tab:heterogeneity}
\begin{threeparttable}
\begin{tabular}{lccc}
\toprule
分组 & 增量AUC & Pseudo R$^2$提升 & 观测数 \\
\midrule
\textbf{Panel A: 按市值规模} & & & \\
小市值（$\le$30\%分位） & 0.085$^{***}$ & 0.042 & 29\,556 \\
中市值（30\%—70\%分位） & 0.062$^{***}$ & 0.031 & 39\,408 \\
大市值（$\ge$70\%分位） & 0.038$^{**}$ & 0.019 & 29\,556 \\
\midrule
\textbf{Panel B: 按分析师覆盖} & & & \\
高覆盖（$\ge$3人） & 0.042$^{**}$ & 0.022 & 32\,512 \\
低覆盖（1—2人） & 0.065$^{***}$ & 0.035 & 31\,218 \\
无覆盖（0人） & 0.091$^{***}$ & 0.048 & 25\,015 \\
\midrule
\textbf{Panel C: 按行业} & & & \\
制造业 & 0.058$^{***}$ & 0.031 & 52\,216 \\
TMT行业 & 0.072$^{***}$ & 0.039 & 18\,719 \\
其他 & 0.048$^{***}$ & 0.025 & 27\,585 \\
\bottomrule
\end{tabular}
\begin{tablenotes}
\item[注] 增量AUC为在Base模型（控制变量+财务指标）基础上加入文本语义特征后的AUC增加值。分组回归模型均包含行业和时间固定效应。$^{***}$ p<0.01, $^{**}$ p<0.05, $^{*}$ p<0.1。
\end{tablenotes}
\end{threeparttable}
\end{table}

这些结果具有直观的经济学解释：信息不对称是软信息价值产生的前提条件。当一家公司被广泛覆盖（如大市值、高分析师关注）时，其已公开的信息已经较为充分，软信息的"增量"空间有限。相反，对于"信息孤岛"式的小公司，公开财务指标往往难以完整反映其风险状况，这时公告文本中蕴含的软信息就成为重要的信息补充渠道。